\documentclass[11pt]{article}
\usepackage{mathblatt}
% vorspann.tex – Kompetenzblatt, zusammenbau v0.7 (2026-09-28).
% Ergänzt mathblatt.sty für Kompetenzblätter, ohne sie zu ändern
% (layout-befunde.md 1–34 vom 28.09.). Wird nach \usepackage{mathblatt}
% eingelesen.
\makeatletter
\RequirePackage{newunicodechar}
% Zeichen, die Latin Modern nicht hat (Befund 20): Text aus einer
% Ersatzschrift, Mathe als Befehl; ohne Ersatzschrift auch Text als Mathe.
\newif\ifkbersatz
\IfFontExistsTF{FreeSerif}{\newfontfamily\kbersatzschrift{FreeSerif}\kbersatztrue}{%
 \IfFontExistsTF{DejaVu Serif}{\newfontfamily\kbersatzschrift{DejaVu Serif}\kbersatztrue}{%
  \IfFontExistsTF{Cambria Math}{\newfontfamily\kbersatzschrift{Cambria Math}\kbersatztrue}{%
   \IfFontExistsTF{Segoe UI Symbol}{\newfontfamily\kbersatzschrift{Segoe UI Symbol}\kbersatztrue}{}}}}
\newcommand{\kbzeichen}[2]{\ifmmode #2\else\ifkbersatz{\kbersatzschrift #1}\else\ensuremath{#2}\fi\fi}
\newcommand{\kbzeichenhoch}[1]{\ifkbersatz\ifmmode\text{\kbersatzschrift #1}\else{\kbersatzschrift #1}\fi\else ?\fi}
% Kopf- und Fußzeile (Befund 1, 2): Kopf Thema · Blattart, Fuß Kennung und Seite
\newcommand{\kbkopfzeile}[2]{\pagestyle{fancy}\fancyhf{}\renewcommand{\headrulewidth}{0pt}%
  \fancyhead[L]{\footnotesize\color{mbgrau}#1}%
  \fancyfoot[L]{\footnotesize\color{mbgrau}#2}%
  \fancyfoot[R]{\footnotesize\color{mbgrau}Seite \thepage}}
% Flattersatz überall (Befund 25)
\raggedright
% Titel des Blatts: Ich-kann-Satz, darunter Zweigzeile
\newcommand{\kbtitel}[2]{\par\noindent{\Large\bfseries\raggedright #1\par}%
  \ifblank{#2}{}{\par\smallskip\noindent{\small #2}\par}\medskip}
% Abschnitt: „Das kennst du schon“, „Zum Merken“
\newcommand{\kbabschnitt}[1]{\par\addvspace{12pt}\Needspace*{5\baselineskip}%
  \noindent{\bfseries #1}\par\nobreak\vspace{3pt}}
% Hauptnummer: Titel hängt am ersten Block; jeder Block (Teilaufgabe) bricht
% nicht, passt er nicht mehr, rückt er ganz auf die nächste Seite (Befund 21).
\newif\ifkbtitel
\newsavebox{\kbbox}
\newlength{\kbhoehe}
\newlength{\kbnummerabstand}\setlength{\kbnummerabstand}{10pt}
\newlength{\kbteilabstand}\setlength{\kbteilabstand}{6pt}
\newenvironment{kbaufgabe}[1]{\stepcounter{aufgabe}\setcounter{teil}{0}%
  \gdef\kbtiteltext{\noindent\hangindent1.6em\hangafter1\makebox[1.6em][l]{\textbf{\theaufgabe.}}#1\par}%
  \global\kbtiteltrue\par\addvspace{\kbnummerabstand}}{\par}
\newenvironment{kbblock}{\par\addvspace{\kbteilabstand}\begin{lrbox}{\kbbox}%
  \begin{minipage}[t]{\linewidth}\raggedright\parskip\z@
  \ifkbtitel\kbtiteltext\vspace{2pt}\global\kbtitelfalse\fi
  \list{}{\leftmargin1.6em\labelwidth1.4em\labelsep0.2em\topsep\z@\partopsep\z@
    \parsep\z@\itemsep\z@\rightmargin\z@}\raggedright}%
  {\endlist\end{minipage}\end{lrbox}%
  \setlength{\kbhoehe}{\dimexpr\ht\kbbox+\dp\kbbox\relax}%
  \Needspace*{\kbhoehe}\noindent\usebox{\kbbox}\par}
% Paarsatz (Platz nutzen, Befund 27): zwei Hauptnummern oder zwei
% Teilaufgaben mit Zeichenaufgabe nebeneinander, als ein Block
\newcommand{\kbnummer}[1]{\stepcounter{aufgabe}\setcounter{teil}{0}%
  \noindent\hangindent1.6em\hangafter1\makebox[1.6em][l]{\textbf{\theaufgabe.}}#1\par\vspace{2pt}}
\newcommand{\kbhalb}[1]{\list{}{\leftmargin1.6em\labelwidth1.4em\labelsep0.2em\topsep\z@
  \partopsep\z@\parsep\z@\itemsep\z@}\raggedright #1\endlist}
\newcommand{\kbpaar}[2]{\par\addvspace{\kbnummerabstand}\begin{lrbox}{\kbbox}%
  \begin{minipage}[t]{0.485\linewidth}\raggedright\parskip\z@ #1\end{minipage}\hfill
  \begin{minipage}[t]{0.485\linewidth}\raggedright\parskip\z@ #2\end{minipage}\end{lrbox}%
  \setlength{\kbhoehe}{\dimexpr\ht\kbbox+\dp\kbbox\relax}\Needspace*{\kbhoehe}\noindent\usebox{\kbbox}\par}
\newcommand{\kbteilpaar}[2]{\item[]\noindent
  \begin{minipage}[t]{0.485\linewidth}\kbhalb{#1}\end{minipage}\hfill
  \begin{minipage}[t]{0.485\linewidth}\kbhalb{#2}\end{minipage}\par}
% Anweisung der Hauptnummer, eigene Zeile über der Teilaufgabe (Befund 5)
\newcommand{\kbanweisung}[1]{\item[]#1\par\vspace{2pt}}
% Kopfzeile einer Teilaufgabe: Buchstabe, Auftakt (halbfett oder Kontext),
% Prüfkennung klein rechts (Befund 3, 12, Beschluss c)
\newlength{\kbkennbreite}\setlength{\kbkennbreite}{1.9cm}
\newcommand{\kbkennung}[1]{{\footnotesize #1}}
\newcommand{\kbteil}[3]{\item[#1]\ifblank{#3}%
  {\parbox[t]{\linewidth}{\raggedright #2}}%
  {\parbox[t]{\dimexpr\linewidth-\kbkennbreite\relax}{\raggedright #2}\hfill
   \parbox[t]{\kbkennbreite}{\raggedleft\kbkennung{#3}}}\par}
% Frage in eigener Zeile (Befund 4), ein Satz je Zeile (Befund 10)
\newcommand{\kbfrage}[1]{\par\vspace{2pt}#1\par}
% Antwortfeld in eigener Zeile darunter, linksbündig (Befund 4, 8)
\newcommand{\kbantwort}[1]{\par\vspace{4pt}\noindent#1\par}
% Zwei Spalten: links Grafik/Tabelle/Aufgabe, rechts Antwort oder Rechenraum
% (Befund 27); oben bündig. Beginnt einen eigenen Listenpunkt ohne Marke.
\newcommand{\kbzweispaltig}[3][0.48]{\item[]\noindent
  \begin{minipage}[t]{#1\linewidth}\raggedright #2\end{minipage}\hfill
  \begin{minipage}[t]{\dimexpr0.98\linewidth-#1\linewidth\relax}\raggedright #3\end{minipage}\par}
% Linke Spalte mit Teilaufgabe (Buchstabe hängt links heraus wie in der Liste)
\newcommand{\kbspalte}[1]{\list{}{\leftmargin\z@\labelwidth1.4em\labelsep0.2em\topsep\z@
  \partopsep\z@\parsep\z@\itemsep\z@}\raggedright #1\endlist}
% Antwortfelder: 2,2 cm, in Punkten 1,2 cm; ohne Umbruch davor
\newcommand{\kbfeld}[1][]{\mbox{\underline{\hspace{2.2cm}}\ifblank{#1}{}{\,#1}}}
\newcommand{\kbfeldk}[1][]{\mbox{\underline{\hspace{1.2cm}}\ifblank{#1}{}{\,#1}}}
% Grafik oder Tabelle unter dem Text, linksbündig (Befund 8, 28)
\newcommand{\kbdarunter}[1]{\par\vspace{4pt}\noindent#1\par}
% Ankreuzen: je Option eine Zeile, linksbündig (Befund 6, 7)
\newcommand{\kbkreuz}[1]{\par\vspace{2pt}\noindent$\square$\ #1\par}
% Bearbeitungsraum (Befund 9): Schreibzeilen wie mathblatt (9 mm)
\newcommand{\kbraum}[1]{\schreibzeilen{#1}}
% Wertetabelle mit Köpfen im Textmodus (Befund 26): wie \wertetabelle der
% Vorlage, aber die Köpfe setzt der Aufruf selbst ($x$, „Zeit in h“).
\NewDocumentCommand{\kbwertetabelle}{o m m m}{%
  \def\mbkopf{}\def\mbwerte{}\def\mbleer{}\setcounter{mbn}{0}%
  \foreach \v in {#4}{\xappto\mbkopf{& $\v$}\xappto\mbleer{& }\stepcounter{mbn}\xdef\mbnn{\arabic{mbn}}}%
  \IfNoValueTF{#1}{}{\foreach \v in {#1}{\ifdefempty{\v}{\xappto\mbwerte{& }}{\xappto\mbwerte{& $\v$}}}}%
  \begingroup\renewcommand{\arraystretch}{1.3}%
  \edef\mbpre{|>{\noexpand\columncolor{mbkasten}}l|*{\mbnn}{>{\noexpand\centering\noexpand\arraybackslash}p{\noexpand\mbzell}|}}%
  \expandafter\mbtabstart\expandafter{\mbpre}\hline
  #2 \mbkopf \\ \hline
  \IfNoValueTF{#1}{\rule{0pt}{\mbzeile}#3 \mbleer \\ \hline}{#3 \mbwerte \\ \hline}
  \end{tabular}\endgroup}
% Merkkasten am Ende (Aufbau des Kompetenzblatts)
\newcommand{\kbmerk}[2]{\par\addvspace{14pt}\begin{lrbox}{\kbbox}\begin{minipage}[t]{\linewidth}%
  \noindent{\bfseries #1}\par\vspace{4pt}%
  \uebersichtskasten{\raggedright #2}\end{minipage}\end{lrbox}%
  \setlength{\kbhoehe}{\dimexpr\ht\kbbox+\dp\kbbox\relax}\Needspace*{\kbhoehe}\noindent\usebox{\kbbox}\par}
% Lösungsblatt: Nummer und Buchstabe halbfett, Lösung daneben
\newcommand{\kbloesung}[2]{\par\addvspace{3pt}\noindent\hangindent2.8em\hangafter1\makebox[2.8em][l]{\textbf{#1}}#2\par}
\makeatother

\begin{document}
\kbkopfzeile{Prozentrechnung · Kompetenzblatt}{PRZ-K1}
% Kompetenzblatt PRZ-K1: prozentrechnung e2 „Prozentsatz“ (zusammenbau v0.7, Niveau FOR)
\kbtitel{Ich kann berechnen, wie viel Prozent ein Teil vom Ganzen ist.}{ab Kl. 7 · P10 ×5}
\setcounter{aufgabe}{0}

\kbabschnitt{Das kennst du schon}

% zone: prozentrechnung-zone-f1-v1
\begin{kbaufgabe}{Ich kann einen Anteil als Bruch schreiben und kürzen.}
\begin{kbblock}
\kbzweispaltig[0.46]{\kbspalte{\kbteil{}{{\bfseries\boldmath $6$ von $10$ Bonbons sind rot}}{}\kbfrage{Welcher Anteil? \\ Kürze.}\kbantwort{\kbfeld}}}{\kbraum{2}}
\end{kbblock}
\end{kbaufgabe}

% zone: prozentrechnung-zone-f3-v1
\begin{kbaufgabe}{Ich kann durch 100 teilen.}
\begin{kbblock}
\kbanweisung{Berechne.}
\kbteil{}{$600 : 100$}{}
\kbantwort{\kbfeld}
\end{kbblock}
\end{kbaufgabe}

% zone: prozentrechnung-zone-f5-v1
\begin{kbaufgabe}{Ich kann auf eine Stelle nach dem Komma runden.}
\begin{kbblock}
\kbanweisung{Runde.}
\kbteil{}{{\bfseries\boldmath $4{,}73$}}{}
\kbantwort{\kbfeld}
\end{kbblock}
\end{kbaufgabe}

\kbabschnitt{Schritt für Schritt}

% leiter: prozentrechnung-e2-k3-s0-v1
\begin{kbaufgabe}{Ich kann den Prozentsatz am Streifen ablesen.}
\begin{kbblock}
\kbteil{}{{\bfseries\boldmath Streifen: $50$ Kinder, grau: $10$ Kinder}}{}
\kbfrage{Wie viel Prozent?}
\kbdarunter{\streifenfeld{20}{$0$}{$50$ Kinder}}
\end{kbblock}
\end{kbaufgabe}

% leiter: prozentrechnung-e2-k3-s1-v1
\begin{kbaufgabe}{Ich kann den Prozentsatz angeben, wenn das Ganze 100 ist.}
\begin{kbblock}
\kbteil{}{{\bfseries\boldmath $17$ von $100$ Schülern spielen ein Instrument}}{}
\kbfrage{Wie viel Prozent?}
\kbantwort{\kbfeld[\%]}
\end{kbblock}
\end{kbaufgabe}

% leiter: prozentrechnung-e2-k3-s2-v1
\begin{kbaufgabe}{Ich kann auf 100 erweitern und den Prozentsatz angeben.}
\begin{kbblock}
\kbzweispaltig[0.46]{\kbspalte{\kbteil{}{{\bfseries\boldmath $19$ von $50$ Kindern}}{}\kbfrage{Wie viel Prozent?}\kbantwort{\kbfeld[\%]}}}{\kbraum{2}}
\end{kbblock}
\end{kbaufgabe}

% leiter: prozentrechnung-e2-k3-s3-v1
\begin{kbaufgabe}{Ich kann den Prozentsatz mit Teil durch Ganzes berechnen.}
\begin{kbblock}
\kbteil{}{{\bfseries\boldmath $27$ von $60$ Minuten Training sind Laufen}}{}
\kbfrage{Wie viel Prozent?}
\kbzweispaltig[0.46]{\kbantwort{\kbfeld[\%]}}{\kbraum{2}}
\end{kbblock}
\end{kbaufgabe}

% leiter: prozentrechnung-e2-k3-s4-v1
\begin{kbaufgabe}{Ich kann den Prozentsatz auf eine Stelle runden.}
\begin{kbblock}
\kbteil{}{{\bfseries\boldmath $17$ von $24$ Kindern waren im Schwimmbad}}{}
\kbfrage{Wie viel Prozent? \\ Runde auf eine Stelle.}
\kbzweispaltig[0.46]{\kbantwort{\kbfeld[\%]}}{\kbraum{2}}
\end{kbblock}
\end{kbaufgabe}

% leiter: prozentrechnung-e2-k3-s5-v1
\begin{kbaufgabe}{Ich kann erst das Ganze bilden und dann den Prozentsatz berechnen.}
\begin{kbblock}
\kbteil{}{{\bfseries\boldmath Im Chor singen $14$ Mädchen und $6$ Jungen}}{}
\kbfrage{Wie viel Prozent sind Jungen?}
\kbzweispaltig[0.46]{\kbantwort{\kbfeld[\%]}}{\kbraum{2}}
\end{kbblock}
\end{kbaufgabe}

% leiter: prozentrechnung-e2-k3-s6-v1
\begin{kbaufgabe}{Ich kann den Rabatt in Prozent berechnen.}
\begin{kbblock}
\kbzweispaltig[0.46]{\kbspalte{\kbteil{}{{\bfseries\boldmath Jacke: vorher $80\,€$, jetzt $60\,€$}}{}\kbfrage{Wie viel Prozent Rabatt?}\kbantwort{\kbfeld[\%]}}}{\kbraum{2}}
\end{kbblock}
\end{kbaufgabe}

% pruefung: prozentrechnung-e2-k3-s9-v1, prozentrechnung-e2-k3-s7-v3, prozentrechnung-e2-k3-s7-v1, prozentrechnung-e2-k3-s7-v5, prozentrechnung-e2-k3-s8-v1
\begin{kbaufgabe}{Ich kann das auch in Aufgaben aus der Prüfung.}
\begin{kbblock}
\kbteil{a)}{Körpergrößen der $13$ Spieler einer Mannschaft in cm: $172$, $185$, $169$, $190$, $178$, $181$, $166$, $188$, $175$, $183$, $170$, $192$, $177$.}{P10 ’25}
\kbfrage{Wie viel Prozent sind größer als $180$ cm? \\ Runde auf eine Stelle.}
\item[]\kbraum{2}
\end{kbblock}
\begin{kbblock}
\kbteil{b)}{{\bfseries\boldmath Einwohner einer Stadt}}{P10 ’23}
\kbfrage{Gesamt $480\,000$; Kinder $72\,000$; Jugendliche $57\,000$; Erwachsene $264\,000$; Senioren $87\,000$. \\ Wie viel Prozent sind Jugendliche? \\ Runde auf eine Stelle.}
\item[]\kbraum{2}
\end{kbblock}
\begin{kbblock}
\kbteil{c)}{Bei einer Tombola gibt es $21$ Nieten und $3$ Hauptgewinne.}{P10 ’18}
\kbfrage{Wie viel Prozent der Lose sind Hauptgewinne?}
\kbzweispaltig[0.46]{\kbantwort{\kbfeld[\%]}}{\kbraum{2}}
\end{kbblock}
\begin{kbblock}
\kbteil{d)}{Von $64$ Bewerbungen wurden $47$ angenommen.}{P10 ’15}
\kbfrage{Wie viel Prozent wurden abgelehnt? \\ Runde auf eine Stelle.}
\item[]\kbraum{2}
\end{kbblock}
\begin{kbblock}
\kbteil{e)}{Eine Stadt hat $40\,000$ Einwohner, $1\,080$ sind in diesem Jahr zugezogen.}{P10 ’14}
\kbfrage{Wie viel Prozent sind das?}
\kbzweispaltig[0.46]{\kbantwort{\kbfeld[\%]}}{\kbraum{2}}
\end{kbblock}
\end{kbaufgabe}

\kbmerk{Zum Merken}{Prozentsatz: Teil geteilt durch Ganzes – das ist der Anteil, dann in Prozent schreiben. \\ 14 von 40: 14 : 40 = 0,35 = 35 \% \quad 3 von 8: 3 : 8 = 0,375 = 37,5 \%}
\end{document}
